% TOP_PROBLEM: {"id": 2114, "title": "Erdős Problem #425", "category_id": 1, "category": {"id": 1, "name": "number_theory", "display_name": "Number Theory", "description": "Properties of integers, prime numbers, Diophantine equations.", "slug": "number-theory", "order_index": 1, "created_at": "2026-07-31T15:26:25.670Z"}, "status": "open"}
% FIRST_POSTED: null
% ENTRY_KIND: statement_only
% Imported from: unsolvedmath_additional_500.tex; UnsolvedMath ID 2114
% !TeX program = lualatex
% Compile twice: lualatex 2114.tex
% Self-contained: no external bibliography, images, or \input files.
% Numeric section labels and visible numbers are original UnsolvedMath IDs.
\documentclass[11pt,a4paper]{article}
\usepackage[margin=24mm,headheight=14pt]{geometry}
\usepackage{fontspec}
\setmainfont{Latin Modern Roman}
\setsansfont{Latin Modern Sans}
\setmonofont{Latin Modern Mono}
\usepackage{amsmath,amssymb,mathtools}
\usepackage{unicode-math}
\setmathfont{Latin Modern Math}
\usepackage{microtype}
\usepackage{enumitem,xurl,needspace}
\usepackage[dvipsnames]{xcolor}
\usepackage{hyperref}
\hypersetup{unicode=true,colorlinks=true,linkcolor=MidnightBlue,urlcolor=MidnightBlue,
 pdftitle={UnsolvedMath Problem 2114},pdfauthor={Source-annotated compilation},bookmarksnumbered=true}
\usepackage{bookmark,tocloft,titlesec,fancyhdr}
\setlength{\cftsecnumwidth}{4.2em}
\cftsetpnumwidth{2.5em}
\cftsetrmarg{4em}
\titleformat{\section}{\Large\bfseries\raggedright}{\thesection}{0.7em}{}
\pagestyle{fancy}\fancyhf{}
\fancyhead[L]{\small\sffamily Additional UnsolvedMath statements}
\fancyhead[R]{\small Original catalogue IDs}
\fancyfoot[C]{\thepage}
\setlength{\parindent}{0pt}
\setlength{\parskip}{5pt plus 1pt}
\setlength{\emergencystretch}{3em}
\setcounter{tocdepth}{1}\setcounter{secnumdepth}{1}
\setlist[itemize]{leftmargin=3em,itemsep=4pt,topsep=4pt,parsep=0pt}
\allowdisplaybreaks
\newcommand{\N}{\mathbb{N}}
\newcommand{\Z}{\mathbb{Z}}
\newcommand{\Q}{\mathbb{Q}}
\newcommand{\R}{\mathbb{R}}
\newcommand{\C}{\mathbb{C}}
\newcommand{\F}{\mathbb{F}}
\newcommand{\A}{\mathbb{A}}
\newcommand{\PP}{\mathbb{P}}
\newcommand{\E}{\mathbb{E}}
\newcommand{\eps}{\varepsilon}
\newcommand{\dd}{\,\mathrm d}
\newcommand{\sref}[2]{\hyperlink{source:#1:#2}{[S#2]}}
\newif\ifshowcatalogue
\showcataloguetrue
\newcommand{\cataloguescope}[1]{\ifshowcatalogue
 \paragraph{Statement recorded in the supplied source.}
 #1\fi}
\begin{document}
% UnsolvedMath ID 2114; source catalogue code EP-425

\setcounter{section}{2113}

\section{Sharp Bounds for Multiplicative Sidon Sets}\label{2114}

{\small\sffamily UnsolvedMath ID 2114 · Number Theory}

\subsection{Definitions and mathematical statement}

\paragraph{Shared notation.}Unless a section says otherwise, $\N=\{1,2,\ldots\}$,
$\Z$ is the ring of integers, $\Q,\R,\C$ are the rational, real, and complex
fields, $[N]=\{1,\ldots,\lfloor N\rfloor\}$, and $\log$ is the natural
logarithm. For real functions of a parameter tending to infinity,
$f\ll g$ and $f=O(g)$ mean $|f|\leq Cg$ eventually for some fixed $C$;
$f\gg g$ reverses this comparison; $f=o(g)$ means $f/g\to0$;
$f\sim g$ means $f/g\to1$; and $f\asymp g$ means both order bounds.
A subscript on $O$ or $\ll$ specifies allowed constant dependence.
The expression $n^{o(1)}$ in an upper bound means growth smaller than
$n^\epsilon$ for each fixed $\epsilon>0$. Local definitions govern any
exception, finite-field convention, or use of the same symbol for another
object. An asymptotic claim is not automatically an effective algorithm.



The pair products use distinct factors and are unordered. The first sentence's $N$ and $n$ refer to the same interval-size variable; the display below uses $n$ consistently. In the higher-product question, all factors within a product are distinct. \sref{2114}{1} \sref{2114}{2}

\paragraph{Statement recorded in the supplied source.}
Let $F(n)$ be the maximum possible size of a subset $A\subseteq\{1,\ldots,N\}$ such that the products $ab$ are distinct for all $a<b$. Is there a constant $c$ such that $ F(n)=\pi(n)+(c+o(1))n^{3/4}(\log n)^{-3/2}? $ If $A\subseteq \{1,\ldots,n\}$ is such that all products $a_1\cdots a_r$ are distinct for $a_1<\cdots <a_r$ then is it true that $ \lvert A\rvert \leq \pi(n)+O(n^{\frac{r+1}{2r}})? $ \sref{2114}{1}

\paragraph{Scope and interpretation.}

The interval-size variable has been made consistent as n; this is a typographical normalization of the displayed definition, not a change of domain. \sref{2114}{1}

\paragraph{Residual task reported by the archive.}

The embedded review (2026-08-17) records: Prove or disprove convergence to a constant coefficient and settle the r-fold upper bound. \sref{2114}{1}

{\small\emph{Transcription note.} Only unambiguous escape/formatting damage and nonmathematical serialization debris were repaired in the displayed source statement.}

\subsection{Short English statement}

What are the lower-order terms in the largest integer sets with unique products of distinct pairs, and what bounds hold for longer products?

\subsection{Sources}

{\small

\begin{itemize}

\item[\hypertarget{source:2114:1}{S1}] ULAM.AI, \emph{UnsolvedMath}, supplied \texttt{problems.json}, record 2114 (EP-425), statement and background. The embedded literature-review date is 2026-08-17; that is the archive’s review date, not a fresh certification. Dataset landing page: \url{https://huggingface.co/datasets/ulamai/UnsolvedMath}.

\item[\hypertarget{source:2114:2}{S2}] T. F. Bloom, \emph{Erdős Problems}, Problem 425, \url{https://www.erdosproblems.com/425}. Reference imported from the supplied record; the live problem page was not independently checked for this compilation.

\item[\hypertarget{source:2114:3}{S3}] Hong Liu and Péter Pál Pach, The number of multiplicative Sidon sets of integers, arXiv:1808.06182 (2018). \url{https://arxiv.org/abs/1808.06182}. Bibliographic information imported from the supplied record.

\item[\hypertarget{source:2114:4}{S4}] Erd\H{o}s, P., On some applications of graph theory to number theoretic problems. Publ. Ramanujan Inst. (1968/69), 131-136. Bibliographic information imported from the supplied record.

\item[\hypertarget{source:2114:5}{S5}] Erd\H{o}s, P., Problems and results on combinatorial number theory. A survey of combinatorial theory (Proc. Internat. Sympos., Colorado State Univ., Fort Collins, Colo., 1971) (1973), 117-138. Bibliographic information imported from the supplied record.

\item[\hypertarget{source:2114:6}{S6}] Erd\H{o}s, Paul, Problems and results on combinatorial number theory. III. Number theory day (Proc. Conf., Rockefeller Univ., New York, 1976) (1977), 43-72. Bibliographic information imported from the supplied record.

\end{itemize}

}

\label{end:2114}
\end{document}
